% ============================================================================
\section{Resolución: del renglón al concepto y a sus candidatos}
\label{sec:resolucion}
% ============================================================================

La resolución decide, para cada renglón que se busca en el catálogo,
\begin{enumerate*}[label=(\alph*)]
  \item a qué concepto corresponde,
  \item qué atributos exige,
  \item si el catálogo puede cumplirlo y
  \item qué productos ofrecer.
\end{enumerate*}
Es la etapa donde el sistema \emph{entiende}, y donde un error es más costoso
de detectar: un renglón mal clasificado produce un producto plausible pero
equivocado. Por eso la resolución, más que ninguna otra etapa, está construida
alrededor de la abstención.

\subsection{Definición del problema}

\begin{definicion}[Resolución de un renglón]
\label{def:resolucion}
\textbf{Entrada:} el texto $x$ de un renglón buscable; el catálogo
$(\Cat,\Prod)$; opcionalmente, el vector semántico $\mathbf{e}(x)\in\R^{d}$
del renglón y un concepto asignado por una persona $c^\ast$.
\textbf{Salida:} una tupla
\[
  \Phi_{\mathrm{res}}(x) = \bigl(\tilde x,\ \omega,\ c,\ s,\ \mathrm{Imp},\ O\bigr)
\]
donde $\tilde x$ es el texto corregido, $\omega\in\{\textsc{reconocido},
\textsc{abstención}, \textsc{hueco}, \textsc{sin\_match}\}$ el veredicto,
$c\in\Cat\cup\{\varnothing\}$ el concepto, $s\in[0,1]$ el puntaje de parecido,
$\mathrm{Imp}$ los atributos impuestos y $O\subseteq\Prod(c)$ las opciones
ofrecidas, con $\card{O}\leq 3$.
\end{definicion}

El veredicto distingue cuatro situaciones que el acudiente debe ver de forma
distinta:
\begin{itemize}
  \item \textsc{sin\_match}: ningún concepto alcanza el umbral de
        reconocimiento; no se entendió qué se pide.
  \item \textsc{hueco}: el concepto se reconoce, pero no tiene productos.
  \item \textsc{abstención}: el concepto se reconoce y tiene productos, pero
        ninguno declara el atributo que el renglón exige. Se prefiere no
        ofrecer nada a ofrecer otro producto: a «block oficio» no se le
        responde con papel origami.
  \item \textsc{reconocido}: hay al menos una opción que cumple.
\end{itemize}

El algoritmo tiene cinco fases: corrección ortográfica, reescrituras de
dominio, clasificación léxico-semántica, extracción de atributos con filtro
duro y selección de una opción por nivel. Las subsecciones siguientes las
desarrollan en ese orden, y la \cref{sec:res-algoritmo} las reúne.

% ---------------------------------------------------------------------------
\subsection{Corrección ortográfica conservadora}
\label{sec:res-ortografia}
% ---------------------------------------------------------------------------

\paragraph{Léxico.} El léxico $\mathcal{L}$ es el conjunto de palabras con las
que el catálogo nombra lo que vende: los tokens alfabéticos de longitud al
menos 4 de los nombres y sinónimos de todos los conceptos, en forma plana, más
sus singulares (\cref{def:singular}). Un sinónimo mal escrito a propósito
(por ejemplo «carperta») también entra: si la lista trae la misma falta, no hay
nada que corregir.

\paragraph{Distancia.} Se usa la distancia de \emph{alineamiento óptimo de
cadenas} (OSA, también llamada Damerau-Levenshtein restringida)
\cite{damerau1964,levenshtein1966}, que cuenta como una operación la
inserción, la eliminación, la sustitución y la transposición de dos
caracteres adyacentes. Para $a=a_1\cdots a_n$ y $b=b_1\cdots b_m$ se calcula
por programación dinámica \cite{wagner1974}:
\[
  d_{i,j} = \min\begin{cases}
    d_{i-1,j}+1 & \text{(eliminación)}\\
    d_{i,j-1}+1 & \text{(inserción)}\\
    d_{i-1,j-1}+\ind{a_i\neq b_j} & \text{(sustitución)}\\
    d_{i-2,j-2}+1 & \text{si } i,j>1,\ a_i=b_{j-1},\ a_{i-1}=b_j \text{ (transposición)}
  \end{cases}
\]
con $d_{i,0}=i$ y $d_{0,j}=j$, y la distancia es $d_{n,m}$. La transposición es
la razón de usar OSA y no Levenshtein puro: «taolla» y «toalla» son el mismo
error de dedo. Con Levenshtein quedan a distancia 2 y la corrección no se
haría; con OSA quedan a distancia 1. Como solo interesa saber si la distancia
es a lo sumo 1, el cálculo se corta: si $\abs{n-m}>1$, la distancia es al
menos 2 y no se llena la tabla.

\begin{definicion}[Corrección única]
\label{def:correccion}
Sea $w$ un token en forma plana. La corrección $\chi(w)$ se define así:
\begin{enumerate}
  \item Si $\card{w}<5$, o $w\in\mathcal{L}$, o $\sing(w)\in\mathcal{L}$:
        $\chi(w)=w$ (no se toca).
  \item $V\gets\{v\in\mathcal{L} : v_1=w_1,\ d_{\mathrm{OSA}}(w,v)\leq1\}$,
        las palabras del léxico a una operación que comparten la primera letra.
  \item $V'\gets\{v\in V : \card{v}=\card{w}\}$. Si $V'=\varnothing$:
        $\chi(w)=w$.
  \item Si $\card{V'}=1$, $\chi(w)$ es el único elemento de $V'$; en otro caso,
        $\chi(w)=w$.
\end{enumerate}
El texto corregido $\tilde x$ reemplaza cada token alfabético $w$ de $x$ por
$\chi(w)$, preservando el uso de mayúsculas del original.
\end{definicion}

\begin{proposicion}[Clase de errores corregibles]
\label{prop:correccion-clase}
Si $\chi(w)=v\neq w$, entonces $v$ se obtiene de $w$ mediante exactamente una
sustitución de un carácter distinto del primero o una transposición de dos
caracteres adyacentes. En particular, $\chi$ nunca corrige inserciones ni
eliminaciones.
\end{proposicion}

\begin{proof}
Por el paso 3, $\card{v}=\card{w}$, y por el paso 2,
$d_{\mathrm{OSA}}(w,v)=1$, porque $d=0$ implicaría $v=w\in\mathcal{L}$, caso
excluido en el paso 1. Una inserción o una eliminación cambian la longitud,
así que la única operación posible es una sustitución o una transposición.
Como $v_1=w_1$, la sustitución no afecta el primer carácter. Una
transposición de las posiciones 1 y 2 cambiaría el primer carácter, salvo que
ambos fueran iguales, en cuyo caso la transposición sería la identidad.
\end{proof}

La definición es deliberadamente estrecha, y cada restricción proviene de un
caso real:
\begin{itemize}
  \item \emph{Unicidad.} Si dos palabras del léxico están a una operación, la
        corrección no se atreve: es preferible no reconocer a reconocer mal.
  \item \emph{Primera letra.} «nacho» está a una operación de «ancho», pero no
        comparten la primera letra. Eso no es un dedo corrido sino otra
        palabra.
  \item \emph{Misma longitud.} «taolla» está a una operación de «toalla»
        (transposición) y de «talla» (eliminación). Quedarse con la de igual
        longitud elige la transposición. Además evita reescribir «paquetes»
        como «paquete», lo que cambiaría la cabeza de la frase cuando el
        clasificador ya singulariza por su cuenta.
\end{itemize}

\paragraph{Complejidad.} Para un token de longitud $\ell$, el filtro por
primera letra y por diferencia de longitud deja un subconjunto
$\mathcal{L}_w\subseteq\mathcal{L}$, y cada distancia cuesta
$O(\ell^2)$. El costo por token es $O(\card{\mathcal{L}}+\card{\mathcal{L}_w}\,\ell^2)$
y el de un renglón con $T$ tokens es $O\bigl(T(\card{\mathcal{L}}+\card{\mathcal{L}_w}\ell^2)\bigr)$.
Con un léxico de unos pocos miles de palabras y $\ell\leq 15$, es
despreciable frente a la clasificación.

% ---------------------------------------------------------------------------
\subsection{Reescrituras de dominio}
\label{sec:res-reescrituras}
% ---------------------------------------------------------------------------

Antes de clasificar, dos patrones se reescriben porque su significado no es
la suma de sus palabras:
\begin{itemize}
  \item \textbf{Rangos y series de durezas.} La expresión regular
        \begin{center}\small\verb!\b\d?[hb]{1,2}(\s*(-|,|y|a|al|hasta)\s*\d?[hb]{1,2}\b)+!\end{center}
        reconoce «2H-6B» o «HB, 2H, 2B, 6B» y los reemplaza por «serie
        durezas», que es como se nombra el concepto \emph{kit de lápices de
        dibujo}. Sin la reescritura, «Lápiz 2H-6B» caía en el concepto
        \emph{lápiz} con la dureza 2H exigida, lo que llevaba a una abstención.
  \item \textbf{Pares de escuadras.} «Escuadras 30° y 45°» pide el par, que se
        vende como juego; se reescribe como «escuadras juego». Una escuadra
        sola sigue siendo una escuadra.
\end{itemize}

% ---------------------------------------------------------------------------
\subsection{Clasificación léxica: la cabeza anclada}
\label{sec:res-lexico}
% ---------------------------------------------------------------------------

El canal léxico se apoya en una observación lingüística sobre el español: en
un sintagma nominal, el núcleo va primero. En «pompones de colores» lo pedido
son pompones, y «de colores» solo los describe. La clasificación no premia que
una palabra del concepto aparezca en cualquier parte del renglón, sino que
aparezca \emph{la cosa pedida}.

\begin{definicion}[Consulta y cabeza]
\label{def:cabeza}
La \emph{consulta} $Q(x)$ se obtiene de la forma plana de $x$ eliminando, en
este orden:
\begin{enumerate*}[label=(\roman*)]
  \item la puntuación,
  \item las fracciones ($\backslash d+/\backslash d+$),
  \item el ruido de empaque y cantidad («unidades», «paquetes», «bolsas»,
        «cajas», «pliegos», «hojas», «metros», «x$N$», \dots, con o sin un
        número delante),
  \item las medidas («$N$ gr», «$N$ ml», «$N$ cm», \dots),
  \item un número inicial suelto,
\end{enumerate*}
y quedándose con las primeras seis palabras de longitud mayor que 1 que no
sean palabras vacías. La \emph{cabeza} $\cab(x)$ es la primera palabra de
$Q(x)$, o la cadena vacía si no hay ninguna.
\end{definicion}

\begin{definicion}[Singular]
\label{def:singular}
Para una palabra $w$ en forma plana,
\[
  \sing(w)=\begin{cases}
    w_{1..\card{w}-3}\cdot\texttt{z} & \text{si } \card{w}>4 \text{ y } w \text{ termina en } \texttt{ces},\\
    w_{1..\card{w}-2} & \text{si } \card{w}>4 \text{ y } w \text{ termina en } \texttt{es},\\
    w & \text{si } w \text{ termina en } \texttt{es} \text{ (y } \card{w}\leq4\text{)},\\
    w_{1..\card{w}-1} & \text{si } \card{w}>3 \text{ y } w \text{ termina en } \texttt{s},\\
    w & \text{en otro caso.}
  \end{cases}
\]
Dos palabras \emph{nombran la misma cosa}, $a\approx b$, si sus singulares
coinciden en los primeros cinco caracteres cuando ambos tienen al menos cinco,
o si son iguales en otro caso.
\end{definicion}

La primera regla existe por un error concreto: «lápices de colores» se
clasificaba como \emph{esfero}. Comparando prefijos de cinco caracteres,
«lápices» da «lapic», que no coincide con «lapiz» (el plural de las palabras
en \emph{-z} es \emph{-ces}), pero sí con «lapicero», sinónimo de esfero. Lo
que fallaba no era la longitud del prefijo sino comparar formas flexionadas.
La tercera regla protege monosílabos como «tres» y «mes».

\begin{definicion}[Palabras significativas y parecido de Jaccard]
\label{def:jaccard}
Sea $\mathcal{V}$ el conjunto de palabras vacías (artículos, preposiciones y
conjunciones). Las \emph{palabras significativas} de un texto son
\[
  \sig(x)=\{\sing(w) : w\in\mathrm{tokens}(\bar x),\ w\notin\mathcal{V},\
           w \text{ no es numérico},\ \card{w}>1\},
\]
y el parecido entre dos textos es el coeficiente de Jaccard
\cite{jaccard1912} sobre esos conjuntos:
\[
  J(x,y)=\frac{\card{\sig(x)\cap\sig(y)}}{\card{\sig(x)\cup\sig(y)}},
  \qquad J(x,y)=0 \text{ si alguno es vacío.}
\]
\end{definicion}

Excluir las palabras vacías y los números cambia resultados de manera
apreciable. «chaquiras de colores» y «colores de madera» comparten «colores»
y «de». Con la preposición, el Jaccard llegaba a $0{,}5$, justo el umbral, y
unas chaquiras entraban al concepto de los lápices de colores. Los números no
describen la cosa, pero inflan el denominador y hunden el parecido de las
frases que traen cantidad.

% ---------------------------------------------------------------------------
\subsection{Canal semántico}
\label{sec:res-semantico}
% ---------------------------------------------------------------------------

El canal léxico no puede reconocer un sinónimo que nadie ha escrito en el
catálogo: «estuche» por «cartuchera». Para eso se agrega un canal semántico
basado en \emph{embeddings}: una función $\mathbf{e}\colon\Sigma^*\to
\mathbb{S}^{d-1}$ que lleva un texto a un vector unitario de dimensión
$d=1024$, de modo que textos con significado parecido quedan cerca en ángulo.
En producción, $\mathbf{e}$ es un modelo de \emph{embeddings} de propósito
general \cite{lee2025gemini}, solicitado explícitamente a 1024 dimensiones,
y cada vector se renormaliza a norma 1.

\paragraph{Vector de producto.} Cada producto se embebe a partir de un texto
canónico que concatena su nombre, el nombre de su concepto y sus atributos
declarados en forma «clave valor». El nombre solo no basta: «Cosido 100 hojas»
y «Argollado 100 hojas» se parecen en letras, pero son cuadernos distintos, y
es el atributo el que los separa. El texto embebido se guarda junto al vector
como \emph{huella}, y un vector se recalcula solo si cambió el texto o el
modelo. Así el recálculo sobre todo el catálogo es idempotente y solo paga
lo que cambió.

\paragraph{Vector de concepto.} El concepto no se embebe directamente: su
vector es el \emph{centroide normalizado} de los vectores de sus productos,
\[
  \bar{\mathbf{e}}(c)=\frac{\sum_{p\in\Prod(c)}\mathbf{e}(p)}
                           {\bigl\lVert\sum_{p\in\Prod(c)}\mathbf{e}(p)\bigr\rVert}.
\]
La elección tiene una justificación precisa:

\begin{proposicion}[El centroide normalizado maximiza el coseno total]
\label{prop:centroide}
Sean $\mathbf{v}_1,\dots,\mathbf{v}_n$ vectores unitarios con
$\mathbf{s}=\sum_i\mathbf{v}_i\neq\mathbf{0}$. Entonces
\[
  \frac{\mathbf{s}}{\lVert\mathbf{s}\rVert}
  = \argmax_{\lVert\mathbf{u}\rVert=1}\ \sum_{i=1}^{n}\cos(\mathbf{u},\mathbf{v}_i).
\]
\end{proposicion}

\begin{proof}
Para $\lVert\mathbf{u}\rVert=1$ y $\lVert\mathbf{v}_i\rVert=1$ se tiene
$\cos(\mathbf{u},\mathbf{v}_i)=\mathbf{u}^\top\mathbf{v}_i$, así que la suma es
$\mathbf{u}^\top\mathbf{s}$. Por la desigualdad de Cauchy-Schwarz,
$\mathbf{u}^\top\mathbf{s}\leq\lVert\mathbf{u}\rVert\lVert\mathbf{s}\rVert
=\lVert\mathbf{s}\rVert$, con igualdad si y solo si $\mathbf{u}$ es paralelo a
$\mathbf{s}$ y en su mismo sentido, es decir,
$\mathbf{u}=\mathbf{s}/\lVert\mathbf{s}\rVert$.
\end{proof}

El vector del concepto es entonces la dirección «más parecida en promedio» a
sus productos. Tiene además una ventaja práctica: reutiliza los vectores que
ya existen y no requiere embeber el concepto por separado. El vector del
\emph{renglón} se calcula en el momento, una sola vez por lista y en un solo
lote (una petición). Si la llamada falla, el canal semántico
\emph{degrada} a $\mathbf{e}(x)=\varnothing$ y la clasificación sigue siendo
exactamente la léxica. La degradación queda registrada.

Sea
\[
  \kappa(x,c)=\max\bigl(0,\ \cos(\mathbf{e}(x),\bar{\mathbf{e}}(c))\bigr)
\]
cuando ambos vectores existen, y $\kappa(x,c)=0$ en otro caso. El recorte en
0 expresa que un coseno negativo no significa «menos parecido que nada»: sin
señal semántica, el canal no debe aportar ni restar.

% ---------------------------------------------------------------------------
\subsection{La función de puntaje híbrida}
\label{sec:res-puntaje}
% ---------------------------------------------------------------------------

Sea $h=\cab(\tilde x)$ la cabeza del renglón ya corregido y reescrito. Cada
concepto $c$ se compara contra cada una de sus \emph{frases candidatas}
$y\in\{\mathrm{nom}(c)\}\cup S(c)$. Para una frase $y$ con palabras
$y^{(1)},\dots,y^{(L)}$ se definen dos predicados:
\[
  \mathrm{Ancla}(y) \iff h\neq\varepsilon \wedge h\approx y^{(1)},
  \qquad
  \mathrm{Nombra}(y) \iff h\neq\varepsilon \wedge \exists i:\ h\approx y^{(i)}.
\]

\begin{definicion}[Puntaje híbrido]
\label{def:puntaje}
Con constantes $\alpha=0{,}6$, $\beta=0{,}4$, $\gamma=0{,}15$,
$\theta=0{,}86$ y $\phi=0{,}55$, el puntaje de la frase $y$ del concepto $c$
para el renglón $x$ es
\[
  \psi(x,c,y)=\begin{cases}
    \min\bigl(1,\ \alpha+\beta\,J(\tilde x,y)+\gamma\,\kappa(x,c)\bigr)
      & \text{si } \mathrm{Ancla}(y),\\[2pt]
    \min\bigl(1,\ J(\tilde x,y)+\gamma\,\kappa(x,c)\bigr)
      & \text{si } \mathrm{Nombra}(y)\wedge\neg\mathrm{Ancla}(y),\\[2pt]
    \phi+\beta\,\dfrac{\kappa(x,c)-\theta}{1-\theta}
      & \text{si } \neg\mathrm{Nombra}(y)\wedge\kappa(x,c)\geq\theta,\\[6pt]
    0 & \text{en otro caso.}
  \end{cases}
\]
La tercera rama se llama \emph{rescate semántico} y toma valores en
$[\phi,\ \phi+\beta]=[0{,}55,\ 0{,}95]$. El puntaje del concepto es
$s(x,c)=\max_{y}\psi(x,c,y)$, y el clasificador devuelve
\[
  (\hat c,\hat s)=\Bigl(\argmax_{c\in\Cat}s(x,c),\ \max_{c\in\Cat}s(x,c)\Bigr),
\]
rompiendo los empates a favor del primer concepto en el orden de recorrido.
\end{definicion}

Las ramas codifican una jerarquía de evidencia:

\begin{enumerate}
  \item \textbf{Ancla.} La frase del concepto \emph{empieza} por la cosa
        pedida. Esto garantiza un piso de $\alpha=0{,}6$, por encima del
        umbral de reconocimiento. El resto de la frase ($\beta J$) decide entre
        los conceptos que aciertan la cabeza: sin este término, «lápiz»
        (esfero) y «lápiz colores» empataban con 1,0 y el empate lo rompía el
        orden de la tabla, que equivale a no decidir. El coseno aporta un
        desempate pequeño ($\gamma=0{,}15$) basado en el significado del resto
        de la frase.
  \item \textbf{Nombra.} La frase menciona la cosa pedida, pero no en primera
        posición. El puntaje es el Jaccard, sin piso.
  \item \textbf{Rescate.} La frase \emph{no} menciona la cosa pedida, pero el
        significado es casi el mismo ($\kappa\geq0{,}86$). El puntaje parte de
        un piso de $0{,}55$, apenas por encima del umbral de reconocimiento, y
        crece linealmente con lo que el coseno supera al umbral. Es la única
        vía para que «estuche» responda a «cartuchera» sin que nadie lo haya
        escrito.
  \item \textbf{Nada.} Si la frase no nombra la cabeza y el significado no es
        casi idéntico, el puntaje es 0. Esto cierra el error de un concepto de
        una sola palabra que obtenía $0{,}5$ contra cualquier frase de dos
        palabras que lo mencionara de pasada («chaquiras de colores» contra
        «colores»). Exigir que la cosa pedida esté nombrada resuelve el caso
        con más claridad que subir el umbral, que solo habría movido el
        problema de lugar.
\end{enumerate}

\begin{proposicion}[El canal semántico nunca resta]
\label{prop:hibrido-lexico}
Sea $\psi_0(x,c,y)$ el puntaje con $\kappa\equiv0$ (clasificación puramente
léxica). Entonces $\psi(x,c,y)\geq\psi_0(x,c,y)$ para toda frase $y$, y en
consecuencia $\hat s\geq\hat s_0$, sea cual sea el modelo de
\emph{embeddings}.
\end{proposicion}

\begin{proof}
Por definición, $\kappa\geq0$. En las ramas Ancla y Nombra,
$\psi=\min(1,\psi_0+\gamma\kappa)$ con $\psi_0\leq1$ y $\gamma\kappa\geq0$. En
la rama de rescate, $\psi\geq\phi>0$,
mientras que $\psi_0=0$, porque con $\kappa\equiv0<\theta$ esa frase cae en la
rama «Nada». En la rama «Nada», $\psi=\psi_0=0$. Como el máximo es monótono,
$\hat s\geq\hat s_0$.
\end{proof}

\paragraph{Cómo se distribuyen los cosenos reales.} Para dimensionar los
umbrales se midió, con los vectores del catálogo (684 productos en 117
conceptos), el coseno de cada producto contra el centroide de su propio
concepto (excluyéndose a sí mismo) y contra los centroides de los demás. Los
nombres de producto son textos algo más largos que un renglón, así que esta
medición es una aproximación de lo que ve el clasificador, no una prueba.

\begin{table}[htbp]
\centering\small
\caption{Cosenos producto--centroide con los vectores reales del catálogo.}
\label{tab:cosenos}
\begin{tabular}{@{}lrr@{}}
\toprule
& Concepto ajeno & Concepto propio \\
\midrule
Pares & \num{79344} & 675 \\
Mínimo & 0,436 & 0,676 \\
Mediana & 0,610 & 0,858 \\
Percentil 99 / 99,9 & 0,751 / 0,807 & -- \\
Pares con coseno $\geq0{,}86$ (umbral de rescate) & 6 & 328 \\
Pares con coseno $\geq0{,}93$ & 0 & 60 \\
Pares con coseno negativo & 0 & 0 \\
\bottomrule
\end{tabular}
\end{table}

Dos lecturas. Primera: el modelo nunca produjo cosenos negativos, así que el
recorte de $\kappa$ no cambia ningún resultado actual; lo que hace es que la
\cref{prop:hibrido-lexico} no dependa del modelo. Segunda: los cosenos entre
conceptos \emph{no relacionados} se concentran alrededor de 0,6 y no bajan de
0,43. El modelo ocupa un cono estrecho del espacio, y por eso un umbral de
rescate alto (0,86) es el que separa «la misma cosa dicha de otro modo» de
«un tema escolar parecido».

\begin{hallazgo}{Observaciones sobre los límites de la función de puntaje}
\begin{enumerate}
  \item \emph{Rescate y revisión (abierto).} La intención del piso
        $\phi=0{,}55$ es que un rescate quede sujeto a revisión humana. Con el
        umbral de autoaceptación $\tau=0{,}75$ (\cref{sec:res-estados}), eso
        vale mientras $\phi+0{,}4(\kappa-\theta)/(1-\theta)<\tau$, es decir,
        mientras $\kappa<\theta+(1-\theta)(\tau-\phi)/0{,}4=0{,}93$. Un rescate
        con coseno de $0{,}93$ o más se acepta automáticamente. Su impacto
        medido es bajo: ningún producto alcanzó 0,93 contra un concepto
        ajeno (\cref{tab:cosenos}), y la lista ya pasa por revisión humana
        antes de llegar a una familia. El arreglo correcto es que la respuesta
        del emparejamiento indique que vino de un rescate y que ese renglón vaya
        siempre a revisión. Como eso cambia el contrato entre servicios, se
        difiere. Bajar el techo del rescate por debajo de $\tau$ acoplaría el
        catálogo a un umbral que se configura en otro servicio.
  \item \emph{Rescate frente a ancla.} Un rescate de un concepto puede superar
        a un ancla de otro concepto cuyo Jaccard sea bajo: con $J=0$ y
        $\kappa=0$ el ancla vale $0{,}6$, y el rescate lo supera si
        $\kappa>0{,}8775$. Es coherente con la idea de que un significado casi
        idéntico pesa más que una coincidencia de cabeza sin más contexto,
        pero conviene saberlo al calibrar $\theta$.
\end{enumerate}
\end{hallazgo}

\paragraph{Complejidad de la clasificación.} Sea $\Lambda=\sum_{c}(1+\card{S(c)})$
el número total de frases candidatas y $\bar L$ su longitud media en palabras.
Para cada concepto el coseno se calcula una vez, en $O(d)$, y se reutiliza
para todas sus frases. Para cada frase, los predicados y el Jaccard cuestan
$O(\bar L+\card{\sig(\tilde x)})$ con conjuntos de dispersión. El total por
renglón es
\[
  T_{\mathrm{clas}} = O\bigl(\card{\Cat}\,d + \Lambda\,(\bar L+\card{\sig(\tilde x)})\bigr),
\]
lineal en el tamaño del catálogo. Con unos 150 conceptos y $d=1024$, el
término del coseno es de unos $1{,}5\times10^5$ productos escalares por
renglón, lo que justifica calcularlo una vez por concepto y no por frase. El
índice vectorial aproximado (HNSW \cite{malkov2018hnsw}) que existe en la base
no se usa en este paso: con unos cientos de conceptos, el recorrido exhaustivo
es exacto y más barato que mantener un índice de conceptos aparte.

% ---------------------------------------------------------------------------
\subsection{Atributos impuestos y filtro duro}
\label{sec:res-atributos}
% ---------------------------------------------------------------------------

Reconocido el concepto $\hat c$, se leen del texto \emph{solo} los atributos
que $\hat c$ declara como discriminantes:
\[
  \mathrm{Imp}(x,\hat c)=\mathrm{Extraer}\bigl(\tilde x,\ D(\hat c)\bigr)\setminus\mathcal{X},
\]
donde $\mathcal{X}$ es un conjunto pequeño de pares atributo-valor que se leen
pero no se exigen. Por ejemplo, $(\cod{tipo},\cod{escolar})$: en una lista
escolar casi todo dice «escolar», y exigirlo dejaba sin opciones a 27 maletas
que no lo declaraban. $\mathrm{Extraer}$ se rige por tres reglas:

\begin{enumerate}
  \item \textbf{Solo lo pedido.} Se devuelven únicamente los atributos de
        $D(\hat c)$. Lo demás sería ruido en la comparación.
  \item \textbf{Ausente antes que inventado.} Si un atributo no se puede leer,
        no aparece. Nunca se asigna «desconocido» ni el valor más frecuente:
        un atributo inventado lleva a elegir con confianza un producto
        equivocado; uno ausente lleva a revisión.
  \item \textbf{Valor único.} Para atributos de una sola dimensión (color,
        forma, rayado, material, punta, presentación, encuadernación, grosor),
        el valor se impone solo si el texto nombra \emph{exactamente uno}. Las
        coincidencias se toman de la más larga a la más corta y una que se
        superpone con otra ya tomada no cuenta («doble línea» es un rayado, no
        «doble línea» más «línea»). Si el texto enumera varios valores, el
        renglón pide un surtido. Exigir el primero sería exigir lo contrario:
        «pinceles planos, redondos y abanico» pedía solo planos. En el caso del
        color, si el texto dice «surtido», «multicolor» o «de colores», se
        impone \cod{color=surtido}.
\end{enumerate}

\begin{definicion}[Cumplimiento]
Un producto $p$ \emph{cumple} los atributos impuestos si
\[
  \mathrm{cumple}(p,\mathrm{Imp}) \iff
  \forall (a,v)\in\mathrm{Imp}:\ \alpha(p)(a)\ \text{está definido y}\
  \mathrm{minúsc}(\alpha(p)(a))=\mathrm{minúsc}(v).
\]
\end{definicion}

El filtro es \emph{duro}: si $\mathrm{Imp}\neq\varnothing$ y ningún producto de
$\Prod(\hat c)$ cumple, el veredicto es \textsc{abstención}, con un motivo
explícito («el renglón pide block con tamaño=oficio y ninguno de los 9
productos lo declara; no ofrezco otro»).

% ---------------------------------------------------------------------------
\subsection{Una opción por nivel}
\label{sec:res-niveles}
% ---------------------------------------------------------------------------

Entre los productos que cumplen se ofrece \emph{uno por nivel}, como máximo
tres. Dentro de cada nivel se prefiere el producto \emph{sin} personaje y, a
igualdad, el de SKU menor:
\begin{gather*}
  O=\Bigl\{\ \min_{\preceq}\{p\in\Prod^{\mathrm{ok}}(\hat c) : \nivel(p)=\ell\}\
  :\ \ell\in\{0,1,2\}\Bigr\},\\
  p\preceq p' \iff \bigl(\ind{\mathrm{per}(p)\neq\varnothing},\mathrm{sku}(p)\bigr)
  \leq_{\mathrm{lex}}\bigl(\ind{\mathrm{per}(p')\neq\varnothing},\mathrm{sku}(p')\bigr),
\end{gather*}
donde $\Prod^{\mathrm{ok}}(\hat c)$ son los productos del concepto que cumplen
los atributos impuestos.
Hay tres razones para limitar a tres opciones. Primero, el acudiente compara
tres cosas, no treinta. Segundo, el optimizador trabaja sobre estas opciones,
y $\card{O}\leq3$ acota el espacio de búsqueda en $3^n$
(\cref{sec:optimizacion}). Tercero, el desempate por SKU hace que la respuesta
no dependa del orden en que la base devolvió las filas. Los productos con
personaje son una vía aparte: no compiten por el puesto de «lo que el colegio
pidió», y se incorporan al kit solo cuando el acudiente pide un personaje
(\cref{sec:opt-personaje}). El precio de cada opción es el menor precio
vigente entre las ofertas activas del producto.

% ---------------------------------------------------------------------------
\subsection{Del veredicto al estado de resolución}
\label{sec:res-estados}
% ---------------------------------------------------------------------------

El veredicto del catálogo se traduce al estado del contrato
(\cref{def:estados}) con un umbral configurable $\tau$, por omisión
$\tau=0{,}75$:
\[
  e(x)=\begin{cases}
    \textsc{ResueltoAuto} & \text{si } \omega=\textsc{reconocido},\ O\neq\varnothing,\ \hat s\geq\tau,\\
    \textsc{Ambiguo} & \text{si } \omega=\textsc{reconocido},\ O\neq\varnothing,\ \hat s<\tau,\\
    \textsc{HuecoCatálogo} & \text{si } \omega\in\{\textsc{hueco},\textsc{abstención}\}
      \text{ o } (\omega=\textsc{reconocido},\ O=\varnothing),\\
    \textsc{SinMatch} & \text{si } \omega=\textsc{sin\_match}.
  \end{cases}
\]
La traducción de \textsc{abstención} a \textsc{HuecoCatálogo} es deliberada:
la abstención significa «te entendí y no tengo lo que pides con ese
atributo», que es un hueco y no una incomprensión.

Los dos umbrales dividen el intervalo de puntajes en tres zonas:
\[
  \underbrace{[0,\ 0{,}5)}_{\textsc{SinMatch}}
  \quad
  \underbrace{[0{,}5,\ \tau)}_{\textsc{Ambiguo}\ \to\ \text{revisión humana}}
  \quad
  \underbrace{[\tau,\ 1]}_{\textsc{ResueltoAuto}}.
\]
Con las constantes vigentes, un ancla sin más palabras en común y sin señal
semántica ($J=0$, $\kappa=0$) vale $0{,}6$ y va a revisión. Para aceptarse
automáticamente, un ancla necesita $J\geq0{,}375$ (sin coseno) o una
combinación equivalente de $J$ y $\kappa$, porque
$0{,}6+0{,}4J+0{,}15\kappa\geq0{,}75\iff 0{,}4J+0{,}15\kappa\geq0{,}15$.

\paragraph{Concepto asignado.} Si una persona ya asignó un concepto $c^\ast$ al
renglón durante la revisión de la lista, el clasificador no se ejecuta:
$\hat c=c^\ast$ y $\hat s=1$. Siguen aplicándose el filtro de atributos y la
selección por nivel. Sin esta regla, la previsualización de los kits ignoraba
la curaduría y decía «no entendimos» a renglones que el revisor ya había
resuelto. El costo es simétrico: si el revisor asigna mal un concepto, el
error también se respeta.

% ---------------------------------------------------------------------------
\subsection{El algoritmo completo}
\label{sec:res-algoritmo}
% ---------------------------------------------------------------------------

\begin{algorithm}[htbp]
\caption{Clasificar un renglón (canal léxico y semántico)}
\label{alg:clasificar}
\begin{algorithmic}[1]
\Require texto corregido $\tilde x$; conceptos $\Cat$ con $\bar{\mathbf{e}}(c)$; vector $\mathbf{e}(x)$ opcional
\Ensure $(\hat c,\hat s)$
\State $\tilde x\gets$ reescribir series de durezas y pares de escuadras
\State $h\gets\cab(\tilde x)$ \Comment{\cref{def:cabeza}}
\State $(\hat c,\hat s)\gets(\varnothing,0)$
\ForAll{$c\in\Cat$}
  \State $\kappa\gets\cos(\mathbf{e}(x),\bar{\mathbf{e}}(c))$ si ambos existen; si no, $0$
  \ForAll{$y\in\{\mathrm{nom}(c)\}\cup S(c)$}
    \State $\mathit{an}\gets h\neq\varepsilon\wedge h\approx y^{(1)}$;\quad
           $\mathit{no}\gets h\neq\varepsilon\wedge\exists i: h\approx y^{(i)}$
    \If{$\mathit{an}$} $\psi\gets\min(1,\ 0{,}6+0{,}4\,J(\tilde x,y)+0{,}15\,\kappa)$
    \ElsIf{$\mathit{no}$} $\psi\gets\min(1,\ J(\tilde x,y)+0{,}15\,\kappa)$
    \ElsIf{$\kappa\geq0{,}86$} $\psi\gets0{,}55+0{,}4\,(\kappa-0{,}86)/0{,}14$
    \Else\ $\psi\gets0$
    \EndIf
    \If{$\psi>\hat s$} $(\hat c,\hat s)\gets(c,\psi)$ \Comment{empate: gana el primero}
    \EndIf
  \EndFor
\EndFor
\State \Return $(\hat c,\hat s)$
\end{algorithmic}
\end{algorithm}

\begin{algorithm}[htbp]
\caption{Emparejar un renglón contra el catálogo}
\label{alg:emparejar}
\begin{algorithmic}[1]
\Require texto $x$; catálogo $(\Cat,\Prod)$; vector $\mathbf{e}(x)$ opcional; concepto asignado $c^\ast$ opcional
\Ensure $(\tilde x,\omega,c,s,\mathrm{Imp},O)$
\State $\tilde x\gets$ reemplazar cada token $w$ de $x$ por $\chi(w)$ \Comment{\cref{def:correccion}}
\If{$c^\ast$ existe en $\Cat$} $(\hat c,\hat s)\gets(c^\ast,1)$
\Else\ $(\hat c,\hat s)\gets\textsc{Clasificar}(\tilde x,\Cat,\mathbf{e}(x))$ \Comment{\cref{alg:clasificar}}
\EndIf
\If{$\hat c=\varnothing$ \OR $\hat s<0{,}5$} \Return $(\tilde x,\textsc{sin\_match},\varnothing,\hat s,\varnothing,\varnothing)$ \EndIf
\State $\mathrm{Imp}\gets\mathrm{Extraer}(\tilde x,D(\hat c))\setminus\mathcal{X}$
\If{$\Prod(\hat c)=\varnothing$} \Return $(\tilde x,\textsc{hueco},\hat c,\hat s,\mathrm{Imp},\varnothing)$ \EndIf
\State $K\gets\Prod(\hat c)$
\If{$\mathrm{Imp}\neq\varnothing$}
  \State $K\gets\{p\in K : \mathrm{cumple}(p,\mathrm{Imp})\}$
  \If{$K=\varnothing$} \Return $(\tilde x,\textsc{abstención},\hat c,\hat s,\mathrm{Imp},\varnothing)$ \EndIf
\EndIf
\State $O\gets$ una opción por nivel de $K$, sin personaje primero, desempate por SKU
\State \Return $(\tilde x,\textsc{reconocido},\hat c,\hat s,\mathrm{Imp},O)$
\end{algorithmic}
\end{algorithm}

\paragraph{Complejidad total.} Para una lista de $m$ renglones buscables, la
resolución cuesta
\[
  T_{\mathrm{res}}(m)=O\Bigl(m\bigl(T_{\mathrm{orto}}+T_{\mathrm{clas}}+\card{\Prod(\hat c)}\,\card{D(\hat c)}\bigr)\Bigr)
  + T_{\mathbf{e}}(m),
\]
donde $T_{\mathbf{e}}(m)$ es una sola llamada por lotes al modelo de
\emph{embeddings}. El índice del catálogo (conceptos, productos, precios
mínimos y centroides) se carga una vez por petición con un número constante de
consultas, en $O(\card{\Prod}\,d)$. Toda la resolución corre contra la base
propia: el código que atiende a un acudiente no tiene ningún camino de red
hacia una tienda, y una prueba automática verifica esa propiedad sobre el
código fuente.

% ---------------------------------------------------------------------------
\subsection{Diseño considerado y no adoptado}
\label{sec:res-alternativas}
% ---------------------------------------------------------------------------

El diseño original de esta etapa proponía una recuperación en tres canales
---léxico por trigramas, denso por vectores y un \emph{grafo} de conceptos con
aristas de sustituibilidad, sinonimia y coocurrencia por grado---, cuya
unión formaría un conjunto de candidatos que después se reordenaría con un
modelo de tipo \emph{cross-encoder}, se calibraría a probabilidad y se
decidiría con dos umbrales. Las mediciones de un prototipo anterior sobre
\num{2710} pares anotados fijaron las líneas base de ese diseño:

\begin{table}[htbp]
\centering\small
\caption{Área bajo la curva ROC por canal en el prototipo anterior.}
\label{tab:auc-prototipo}
\begin{tabular}{@{}lc@{}}
\toprule
Canal & AUC \\
\midrule
Léxico & 0,728 \\
Denso (vectores) & 0,825 \\
Fusión por rango recíproco (RRF) & 0,753 \\
Reordenador \emph{cross-encoder} & 0,879 \\
\bottomrule
\end{tabular}
\end{table}

Dos hallazgos de esa medición determinaron el diseño vigente. El primero es
que la fusión por rango recíproco \cite{cormack2009rrf} medía \emph{peor} que
su propio canal denso, así que no podía ser lo que puntuara. El segundo es que
el 27,7\,\% de los ítems no tenía ningún candidato correcto entre sus cinco
primeros. Eso es \emph{recall} que nunca ocurrió, y ninguna mejora de
ordenamiento lo recupera.

La versión vigente adopta una \emph{fusión por guardas} más conservadora que
la de la \cref{def:puntaje}. La estructura léxica (la cabeza anclada) decide
qué conceptos son admisibles, y el canal semántico solo desempata entre ellos
o rescata un sinónimo con un umbral alto. El grafo de sustituibilidad y el
reordenador no forman parte de la versión vigente: la tabla de aristas existe
en el esquema, pero está vacía, y su incorporación queda como trabajo futuro
(\cref{sec:discusion}).

\begin{evolucion}[: resolución]
\small
\begin{tabular}{@{}L{3.3cm}L{3.9cm}L{5.7cm}@{}}
\toprule
Renglón o síntoma & Causa & Cambio \\
\midrule
«pompones / chaquiras de colores» $\to$ lápices de colores &
la contención de «colores» en cualquier posición puntuaba &
\textbf{cabeza anclada} (\cref{def:cabeza}) y rama «Nada» \\
«lápices de colores» $\to$ esfero &
comparar prefijos de formas flexionadas («lapic») &
\textbf{singularización} (\cref{def:singular}) \\
«2 lápices negros» con parecido 0 &
«lápices» y «lápiz» no compartían palabra &
$\sig(\cdot)$ en singular y sin números \\
«chaquiras de colores» con Jaccard 0,5 &
la preposición «de» contaba como coincidencia &
palabras vacías fuera del Jaccard \\
«metro de cable estañado» $\to$ lentejuelas &
cabeza «metro» &
«metro» como ruido de unidad \\
«(1/8 cartulina)» sin concepto &
cabeza «1/8» &
fracciones fuera de la cabeza; «1/8» se lee como tamaño octavo \\
«Lápiz 2H-6B» en abstención &
se exigía dureza 2H &
reescritura a «serie durezas» \\
«Escuadras 30° y 45°» $\to$ escuadra suelta &
los números no cuentan en el parecido &
reescritura a «escuadras juego» \\
«Maleta escolar» sin opciones &
se exigía \cod{tipo=escolar} &
conjunto $\mathcal{X}$ de valores no exigibles \\
«pinceles planos, redondos y abanico» pedía solo planos &
se tomaba el primer valor &
regla del valor único \\
renglón curado decía «no entendimos» &
se volvía a clasificar el texto &
concepto asignado $c^\ast$ con $\hat s=1$ \\
canal semántico con coseno siempre 0 en producción &
dos copias de la carga del índice; la del servicio no leía vectores &
una sola carga del índice para todos los caminos \\
degradación invisible a solo léxico &
excepción capturada en silencio &
la degradación se registra \\
garantía «híbrido $\geq$ léxico» dependía del modelo &
un coseno negativo restaba en las ramas léxicas &
$\kappa$ recortado en 0 (\cref{prop:hibrido-lexico}) \\
\bottomrule
\end{tabular}

\medskip
Los casos se resolvieron con un criterio explícito: se corrige lo que es
lógico para el dominio, no lo que hace que una prueba deje de fallar. Varias
veces el arreglo cómodo era agregar un sinónimo al concepto que ya tenía
productos para que el kit «saliera completo». Se descartó siempre que el
producto ofrecido no fuera lo que el colegio pedía: una línea que dice «no lo
tenemos» es honesta; un producto equivocado en el carrito no lo es.
\end{evolucion}
